% ===================================================================
%  TABELUL Nr. 13 — Hotelul. — Restaurantul. — Cafeneaua.
%  Traduction 2026 d'apres texto/io, controlee sur texto/fr, texto/en
%  et texto/es. Consigne : voir texto/ro/10-tabelul-01.tex.
%
%  Deux notes, marquees toutes deux « (*) », et deux appels : celui du
%  bloc 01-2 (la chambre) et celui du bloc 09-1 (le menu).
%
%  QUATRE METIERS VOISINS ET QUATRE MOTS DISTINCTS : le conducteur de
%  l'omnibus (30) est « vizitiu », le cocher du fiacre (92)
%  « birjar » ; le chasseur de l'hotel (19) est « comisionar », l'homme
%  de peine du trottoir (39) « hamal ». objekti.py releve les objets
%  par leur nom.
% ===================================================================

%%K t13-noto-1 noto
\VUnotes{143.2mm}{%
\textit{\VUgras{(*) Amănuntele camerei de culcare vezi în tabelul nr. 6.}}%
}

%%K t13-apar-1 apar
\VUsaut{3.0mm}
\VUtitre{15.0pt}{60}{TABELUL Nr. 13}
\VUsaut{1.6mm}
\VUpk{10.2pt}{40}{Hotelul. --- Restaurantul.}
\VUsaut{2.0mm}
\VUpk{10.2pt}{40}{Cafeneaua.}
\VUsaut{3.4mm}

%%K t13-01-1 p
1. --- Ajungând ieri seară la Royan, am tras la \textit{Hotelul Oceanului}, pe care
mi l-a lăudat un prieten.

%%K t13-01-2 p
Mi-au dat o \VUgras{cameră} frumoasă \textsuperscript{(2)} la \VUgras{etajul}
al doilea \textsuperscript{(1)}, unde am dormit foarte bine (*).

%%K t13-02-1 p
2. --- Azi-dimineață, ieșind din camera mea, unde \VUgras{camerista}
\textsuperscript{(3)} îmi adusese gustarea, am intrat în \VUgras{salonul de fumat}
\textsuperscript{(5)}, care slujește și de \VUgras{sală de lectură}
\textsuperscript{(4)}, ca să răsfoiesc \VUgras{ziarele} \textsuperscript{(6)},
fumând o \VUgras{țigară} \textsuperscript{(8)}. După ce am sfârșit de citit, am
coborât cu \VUgras{ascensorul} \textsuperscript{(9)} și am mers să mă plimb prin
oraș.

%%K t13-03-1 p
3. --- Fiindcă uitasem să îl întreb pe \VUgras{funcționarul} hotelului
\textsuperscript{(12)} la ce ceas se servește la \VUgras{masa comună}
\textsuperscript{(11)}, m-am întors devreme și m-am folosit de aceasta ca să fac o
baie în \VUgras{baia} hotelului \textsuperscript{(10)}. Apoi am mers la
\VUgras{casă} \textsuperscript{(15)}, ca să telefonez unui prieten. Dar
\VUgras{telefonul} \textsuperscript{(13)} era ocupat; văzându-mi încurcătura,
\VUgras{casierița} \textsuperscript{(14)}, care trecea o \VUgras{notă}
\textsuperscript{(17)} în \VUgras{registrul călătorilor} \textsuperscript{(16)}, l-a
rugat pe \VUgras{portar} \textsuperscript{(18)} să trimită un
\VUgras{comisionar} \textsuperscript{(19)} să îmi facă treaba pe
\VUgras{bicicleta} lui \textsuperscript{(20)}.

%%K t13-04-1 p
4. --- I-am mulțumit și am ieșit ca să dau un bacșiș \VUgras{hamalului}
\textsuperscript{(24)} (omul bun la toate), care adusese de la gară pe
\VUgras{căruciorul} său de mână \textsuperscript{(25)} \VUgras{cutia mea de
pălării} \textsuperscript{(26)}, \VUgras{geamantanul meu} \textsuperscript{(27)}
și \VUgras{sacul meu de călătorie} \textsuperscript{(28)}. \VUgras{Poștașul}
\textsuperscript{(21)}, care tocmai sosise, mi-a dat o \VUgras{scrisoare}
\textsuperscript{(22)}.

%%K t13-05-1 p
5. --- Am mai zăbovit puțin în prag, lângă \VUgras{cutia de scrisori}
\textsuperscript{(23)}, privind mișcarea de pe stradă. \VUgras{Vizitiul}
\textsuperscript{(30)} \VUgras{omnibuzului} \textsuperscript{(29)} încărca pe
trăsura lui \VUgras{lada} \textsuperscript{(31)} unui \VUgras{călător}
\textsuperscript{(32)}, de care își lua rămas bun \VUgras{stăpânul hotelului}
\textsuperscript{(33)}. Un tânăr \VUgras{telegrafist} \textsuperscript{(35)} ducea,
fără nicio grabă, o \VUgras{depeșă} \textsuperscript{(34)} unuia dintre oaspeții
hotelului; un \VUgras{turist} \textsuperscript{(36)} cu \VUgras{aparatul de
fotografiat} \textsuperscript{(38)} pe umăr își cerceta \VUgras{ghidul}
\textsuperscript{(37)}.

%%K t13-tit-1 sub
\VUsaut{4.0mm}
\VUpk{10.2pt}{40}{Ceasul gustării.}
\VUsaut{3.0mm}

%%K t13-06-1 p
6. --- Înaintea grilajului un \VUgras{hamal} fără grijă \textsuperscript{(39)}
moțăia între \VUgras{cârligele} sale \textsuperscript{(40)} și \VUgras{lada de
lustruit} \textsuperscript{(41)}, în vreme ce o tânără \VUgras{spălătoreasă}
\textsuperscript{(42)}, care ducea un \VUgras{coș} \textsuperscript{(43)} de rufe,
râdea de înfățișarea gravă a \VUgras{șefului de sală} \textsuperscript{(44)}, care
stătea la ușă.

%%K t13-07-1 p
7. --- Vederea \VUgras{bucătarului} \textsuperscript{(45)}, care ieșea din
\VUgras{bucătăria} sa \textsuperscript{(47)} de la \VUgras{subsol}
\textsuperscript{(48)} ca să trimită un \VUgras{ucenic de bucătărie}
\textsuperscript{(46)} cu o treabă, mi-a adus aminte că ceasul gustării este
aproape, și am intrat în restaurant, trecând prin curtea în care un
\VUgras{automobilist} \textsuperscript{(50)} își băga \VUgras{automobilul}
\textsuperscript{(49)}, iar un \VUgras{mecanic} \textsuperscript{(52)} dregea, după
îndrumările \VUgras{biciclistului} \textsuperscript{(53)}, un \VUgras{triciclu cu
benzină} \textsuperscript{(51)}, care tocmai avusese o pană.

%%K t13-08-1 p
8. --- Am întrebat un tânăr \VUgras{grum} \textsuperscript{(54)}, care ducea
\VUgras{halbe de bere} \textsuperscript{(55)}, dacă masa comună este sub verandă,
și, la răspunsul lui că da, m-am așezat la una dintre mese și mi s-a dat o gustare
minunată.

%%K t13-noto-2 noto
\VUnotes{143.2mm}{%
\textit{\VUgras{(*) Cu ajutorul listei de bucate date în vocabular, învățătorul va
putea schimba lesne meniul de mai jos.}}%
}

%%K t13-tit-2 sub
\VUsaut{4.0mm}
\VUpk{10.2pt}{40}{O gustare aleasă.}
\VUsaut{3.0mm}

%%K t13-09-1 p
9. --- Chelnerul mi-a adus meniul gustării (*) și lista vinurilor. Am ales și am
cerut ca \VUgras{aperitiv creveți} cu \VUgras{unt}, iar ca fel întâi
\VUgras{burtă după chipul din Caen}. \VUgras{Pește} nu am luat, dar am cerut o
porție de \VUgras{pulpă} de miel și, ca \VUgras{legumă, spanac} cu smântână. Apoi,
fiindcă niciunul dintre \VUgras{dulciuri} nu mi-a plăcut, am pus să mi se aducă
deserturi felurite: \VUgras{brânză}, \VUgras{fructe} și \VUgras{pesmeți}. Am mai
adăugat un \VUgras{vin} ales de Saint-Émilion și, foarte mulțumit de gustarea mea,
m-am ridicat de la masă după ce i-am dat \VUgras{chelnerului}
\textsuperscript{(57)} un \VUgras{bacșiș} frumos \textsuperscript{(59)}.

%%K t13-10-1 p
10. --- Văzând că mă pregătesc să plec, \VUgras{directorul} \textsuperscript{(60)}
mi-a propus să ies pe balcon ca să admir minunata panoramă a \VUgras{oceanului}
\textsuperscript{(62)}. În depărtare se deosebeau \VUgras{bărci} mici de pescuit
\textsuperscript{(63)}, \VUgras{veliere} \textsuperscript{(64)} și un
\VUgras{pachebot} uriaș \textsuperscript{(65)}, a cărui siluetă neagră se taie pe
\VUgras{norii} albi \textsuperscript{(67)} de la \VUgras{zare}
\textsuperscript{(66)}. Un vânticel ușor mișca \VUgras{steagul}
\textsuperscript{(70)}, înfipt deasupra \VUgras{firmei} \textsuperscript{(69)} a
\VUgras{holului} \textsuperscript{(68)}.

%%K t13-tit-3 sub
\VUsaut{3.0mm}
\VUpk{10.2pt}{40}{Petreceri.}
\VUsaut{3.0mm}

%%K t13-11-1 p
11. --- Priveliștea era atât de plăcută, încât m-am hotărât ca a doua zi să urc pe
terasa cafenelei cu un \VUgras{binoclu de teatru} \textsuperscript{(71)} sau cu o
\VUgras{lunetă} \textsuperscript{(72)}.

%%K t13-12-1 p
12. --- Plecând de pe balcon, am mers la cafeneaua hotelului ca să beau
\VUgras{ceva} \textsuperscript{(73)} și să joc o partidă de \VUgras{biliard}
\textsuperscript{(75)}. Am trecut de-a dreptul prin sala cafenelei, unde mulți
mușterii jucau \VUgras{șah} \textsuperscript{(78)}, \VUgras{dame}
\textsuperscript{(79)}, \VUgras{domino} \textsuperscript{(80)},
\VUgras{table} \textsuperscript{(81)} sau \VUgras{cărți} \textsuperscript{(82)},
și am urcat de-a dreptul în \VUgras{sala de biliard} \textsuperscript{(74)}.

%%K t13-13-1 p
13. --- După ce am sfârșit partida, am coborât la parter și i-am cerut chelnerului
un \VUgras{plic} \textsuperscript{(84)}, \VUgras{hârtie} \textsuperscript{(83)}, un
\VUgras{timbru} \textsuperscript{(85)}, o \VUgras{peniță} \textsuperscript{(87)} și
\VUgras{cerneală} \textsuperscript{(86)}, ca să scriu o scrisoare.

%%K t13-13-2 p
Apoi am chemat un \VUgras{birjar} \textsuperscript{(92)}, a cărui \VUgras{trăsură}
\textsuperscript{(88)} se oprise înaintea cafenelei. L-am întrebat de prețul pe
ceas sau pe cursă, iar când ne-am înțeles la preț, mi-a deschis \VUgras{portiera}
\textsuperscript{(89)} și m-a așezat înăuntru. S-a suit din nou pe \VUgras{capra}
sa \textsuperscript{(91)}, a pornit sprinten caii cu \VUgras{biciușca} sa
\textsuperscript{(93)}, și am plecat într-o plimbare încântătoare.
\VUsaut{6.0mm}
\VUfilet{22mm}
